\documentclass[11pt]{article}
\usepackage[a4paper,margin=1in]{geometry}
\usepackage{fontspec}
\setmainfont{DejaVu Serif}
\usepackage{polyglossia}
\setmainlanguage{english}
\setotherlanguage{french}
\usepackage{amsmath,amssymb,amsthm,mathtools}
\usepackage{tikz-cd}
\usepackage{hyperref}
\hypersetup{colorlinks=true,linkcolor=blue,citecolor=blue,urlcolor=blue}
\newtheorem{theorem}{Theorem}
\newtheorem{corollary}{Corollary}
\newtheorem{proposition}{Proposition}[section]
\newtheorem{lemma}[proposition]{Lemma}
\newtheorem{keylemma}[proposition]{Key Lemma}
\theoremstyle{remark}
\newtheorem{remark}[proposition]{Remark}
\DeclareMathOperator{\Spec}{Spec}
\DeclareMathOperator{\Spf}{Spf}
\DeclareMathOperator{\Lie}{Lie}
\DeclareMathOperator{\Pic}{Pic}
\DeclareMathOperator{\Ext}{Ext}
\DeclareMathOperator{\Hom}{Hom}
\DeclareMathOperator{\Res}{Res}
\newcommand{\OO}{\mathcal O}
\newcommand{\ww}{\omega}
\newcommand{\C}{\mathbf C}
\newcommand{\A}{\mathbf A}
\newcommand{\Gm}{\mathbf G_m}
\newcommand{\GG}{\mathbb G}
\newcommand{\kk}{k}
\DeclareMathOperator{\Tr}{Tr}
\begin{document}

\begin{center}
 {\Large Integration over a Vanishing Cycle}\par
 \vspace{0.5em}
 {P. Deligne}\par
\end{center}

\section*{Introduction}

Let $k$ be a commutative field and let $f$ be a formal power series in two indeterminates $x$ and $y$ over $k$:
\[
 f=\sum a_{n,m}x^ny^m.
\]
Let $I(f)$ be the formal power series in one indeterminate $t$
\[
 I(f)=\sum a_{n,n}t^n.
\]
The aim is to prove the following result.

\begin{theorem}
If $f$ represents an algebraic function of $x$ and $y$, and if $k$ has characteristic $p\ne0$, then $I(f)$ represents an algebraic function of $t$.
\end{theorem}

For a formal series in $N$ indeterminates
\[
 g=\sum a_nx^n,
\]
one defines similarly
\[
 I_N(g)=\sum a_{n,\ldots,n}t^n,
\]
and one has the following consequence.

\begin{corollary}
If $g$ is algebraic over the field $k(x_1,\ldots,x_N)$, and if $k$ has characteristic $p>0$, then $I_N(g)$ is algebraic over $k(t)$.
\end{corollary}

Under the stronger hypothesis $g\in k(x_1,\ldots,x_N)$ this result is due to H. Furstenberg. The first purpose here is to give a geometric explanation of Furstenberg's result.

Suppose first that $k=\mathbb C$ and that the formal series $f$ is convergent. Then $I(f)$ is convergent and has the integral representation
\begin{equation}\tag{1}
 I(f)(t)=\frac1{2\pi i}\int_{xy=t} f(x,y)\,\frac{dx\wedge dy}{dt}.
\end{equation}
Here $\pi:\mathbb C^2\to\mathbb C$ is $(x,y)\mapsto xy$. The expression $(dx\wedge dy)/dt$ is a relative holomorphic differential on $\mathbb C^2-\{(0,0)\}$, relative to $\pi$; it is also a section of the relative dualizing sheaf at $(0,0)$, in the sense of Grothendieck duality. In the coordinate $y$ on the fibre $xy=t$ one has
\[
 \frac{dx\wedge dy}{dt}=\frac{d(t/y)\wedge dy}{dt}=\frac{dy}{y},
\]
up to the sign fixed by the chosen order of the variables.

If $f$ converges in a polydisc of radius $R$, and if $|t|<R^2$, the cycle of integration is $|y|=r$, $|x|=|t|/r$, $xy=t$, with $|t|/R<r<R$. The integral is independent of $r$. Thus
\[
 \frac1{2\pi i}\int_{xy=t} f(x,y)\frac{dx\wedge dy}{dt}
 =\frac1{2\pi i}\int f(t/y,y)\frac{dy}{y},
\]
and expansion into the series followed by the residue formula leaves only the terms with equal exponents in $x$ and $y$.

The integration cycle is the Lefschetz vanishing cycle for $\pi$ near $(0,0)$: if $r=|t|^{1/2}$ the cycle visibly collapses as $|t|\to0$. This explains the title.

Formula (1) suggests that the operation $I$ has an intrinsic meaning for a morphism of formal schemes $\pi:X\to S$ over $k$ with $X\simeq\Spf k[[x,y]]$, $S\simeq\Spf k[[t]]$, and $\pi$ having a non-degenerate quadratic singularity. The two branches of $\pi^{-1}(0)$ are assumed to be defined over $k$, and an order between them is chosen. Thus we choose coordinates with $\pi(x,y)=xy$, with the first branch of $xy=0$ being $x=0$.

Let $\Omega^1_{X/k}$ be the Kähler differentials defined with respect to the adic topology. In coordinates $\Omega^1_{X/k}$ is freely generated by $dx,dy$, and $\Omega^1_{S/k}$ by $dt$. Let
\[
 \omega_{X/k}=\det\Omega^1_{X/k},\qquad
 \omega_{S/k}=\Omega^1_{S/k},\qquad
 \omega_{X/S}=\omega_{X/k}\otimes \pi^*\omega_{S/k}^{-1}.
\]
In coordinates $\omega_{X/S}$ is generated by $(dx\wedge dy)/dt$. The operation $I$ will be interpreted as a morphism
\[
 \pi_*\omega_{X/S}\longrightarrow\OO_S,
\]
or, for $S=\Spf(V)$, as a $V$-linear map
\[
 H^0(X,\omega_{X/S})\longrightarrow V.
\]

Section 1 provides this interpretation algebraically. The integration sign in (1) is replaced by Grothendieck's local trace. For a uniformizer $t$ of $S$, the reduction modulo $t^n$ is denoted with the index $n$. Let $A_n$ be the complete local ring whose formal spectrum is the reduction $X_n$; in coordinates $A_n=k[[x,y]]/((xy)^n)$. It is convenient to replace $X_n$ by the ordinary scheme $X'_n=\Spec A_n$. Let $O$ be its closed point. Grothendieck's local trace is a morphism
\[
 \Tr:H^1_{\{O\}}(X'_n,\omega)\longrightarrow H^0(S_n,\OO).
\]
The role of the vanishing cycle is played by an element
\[
 v_n\in H^1_{\{O\}}(X'_n,\OO).
\]
In coordinates $v_n$ comes from $\Ext^1_{A_n}(A_n/(x+y)^{2n-1},A_n)$. In Section 2 these classes are related to an element of the Lie algebra of a Picard scheme. Section 3 globalizes the construction and proves the theorem and the corollary.

\section{Algebraic analogue of the integral formula}

Keep the notation above. The inclusion $X'_1\subset X'_n$ is a homeomorphism, since $X'_1$ is defined by a nilpotent ideal in $X'_n$. The complement of the origin in $X'_n$ has two connected components corresponding to the branches $x=0$ and $y=0$ of $X'_1$. Let $u_n$ be the section of the structure sheaf of $X'_n-\{O\}$ which is $1$ on the first branch and $0$ on the second.

Algebraically, the homomorphism
\[
 A_n=k[[x,y]]/((xy)^n)	o k[[x,y]]/(x^n)\times k[[x,y]]/(y^n)
\]
induces an isomorphism after inverting $x+y$, and $u_n$ is the element whose image is $(1,0)$. Let $v_n$ be the image of $u_n$ under the boundary map
\[
 H^0(X'_n-\{O\},\OO)\longrightarrow H^1_{\{O\}}(X'_n,\OO).
\]
Applying $\Ext^\bullet(-,A_n)$ to
\[
 0\to ((x+y)^N)\to A_n\to A_n/((x+y)^N)\to0
\]
gives a diagram showing that, for $N=2n-1$, $v_n$ lies in the image of
\[
 \Ext^1_{A_n}(A_n/((x+y)^{2n-1}),A_n).
\]

\paragraph{1.2. Proposition.}
For $f\in k[[x,y]]$ and for every $n$,
\begin{equation}\tag{1.2.1}
 \Tr\left(f\frac{dx\wedge dy}{dt},v_n\right)=-I(f)\pmod {t^n}.
\end{equation}

The trace is Grothendieck's local trace. We prove the formula when $k$ is infinite; the general case follows by extension of scalars. Let $I'(f)$ denote the left-hand side of (1.2.1). It is a continuous $k[[t]]/(t^n)$-linear map $A_n\to k[[t]]/(t^n)$. For $\lambda,\mu\in k^*$ let $\sigma(\lambda,\mu)$ be the automorphism
\[
 (x,y,t)\mapsto(\lambda x,\mu y,\lambda\mu t).
\]
It preserves $(dx\wedge dy)/dt$ and $v_n$, hence
\begin{equation}\tag{1.2.2}
 I'(\sigma(\lambda,mumu)^*f)=\sigma(\lambda,\mu)^*I'(f).
\end{equation}
Applying this to monomials gives $I'(x^ay^b)=0$ unless $a=b$. By continuity,
\[
 I'(f)=I(f)I'(1).
\]
The scalar $I'(1)$ is independent of $n$, so it remains to compute it for $n=1$.

For $n=1$, $X'_1=\Spec k[[x,y]]/(xy)$ is the union of two smooth formal arcs meeting transversely. Sections of $\omega$ are pairs of meromorphic differentials $(\alpha,\beta)$ on the two branches with at most simple poles at $O$ and residues summing to zero. On the punctured curve arbitrary pole orders are allowed. The composite
\[
 H^0(X'_1-\{O\},\omega)\to H^1_{\{O\}}(X'_1,\omega)\xrightarrow{\Tr} k
\]
is the negative of $(\alpha,\beta)\mapsto \Res\alpha+\Res\beta$. In these terms $(dx\wedge dy)/dt$ and $u_1$ give residue sum $1$, proving (1.2.1).

\section{Relation with the Picard functor}

Let $S=\Spec k[[t]]$, let $s$ be its closed point and $\eta$ its generic point. Let $\pi:Y\to S$ be proper and flat, purely of relative dimension one, and let $O\in Y_s(k)$. Assume that the completion of $Y$ at $O$ is the formal singularity considered above. Thus $Y$ is regular at $O$, $O$ is an isolated non-smooth point of $\pi$, the singularity is non-degenerate quadratic, and the two branches are defined over $k$ and ordered.

We may further suppose $Y$ regular and connected. Let $Y_\eta$ be the generic fibre over $k((t))$.

\paragraph{2.2. Lemma.}
The curve $Y_\eta$ is absolutely irreducible and generically smooth over $k((t))$.

Indeed, the special fibre is connected and has a rational point, hence is absolutely connected. It is smooth near $O$ away from $O$, so after a separable extension it has smooth rational points. The corresponding base change has a section; regularity then implies irreducibility, and absolute connectedness gives absolute irreducibility. The smooth locus of the generic fibre is non-empty and therefore dense.

Let $C$ be a complete reduced generically smooth curve over a field $K$, and $D$ a Cartier divisor on $C$. The generalized Jacobian $\Pic(C;D)$ represents invertible sheaves on $C_T$ trivialized on $D_T$. It is smooth, and
\begin{equation}\tag{2.3.1}
 \Lie\Pic(C;D)=H^1(C,\OO(-D)).
\end{equation}
If $C$ is smooth, the map
\begin{equation}\tag{2.3.2}
 j:C-D\to\Pic(C;D),\qquad u\mapsto[\OO(u(T))]
\end{equation}
relates Serre duality with the pullback of invariant differentials on the Picard scheme.

Returning to $Y/S$, let $Z$ be a relative Cartier divisor disjoint from $O$, and let $\Pic(Y;Z)^0$ be the relative generalized Picard scheme parametrizing line bundles on $Y$ trivialized along $Z$ and of degree zero on each geometrically irreducible component of each fibre. For each $n$ one constructs morphisms
\begin{align}\tag{2.4.6}
 H^0(O_T,\OO_T)&\longrightarrow H^1(Y_{n,T},\OO(-Z_{n,T})),\\
\tag{2.4.7}
 H^0(O_T,\OO_T^*)&\longrightarrow H^1(Y_{n,T},\OO(1\text{ on }Z_{n,T})^*).
\end{align}
The multiplicative version gives a morphism of $k[[t]]/(t^n)$-group schemes
\begin{equation}\tag{2.4.8}
 w_n:\Gm\longrightarrow\Pic(Y_n,Z)^0.
\end{equation}
For $m<n$, the reduction of $w_n$ modulo $t^m$ is $w_m$; for $n=1$ this is the morphism associated with the two branches.

Although one cannot take a direct limit to get a morphism $\Gm\to\Pic(Y;Z)^0$, one can restrict to finite subgroup schemes. For every integer $m$ one obtains
\begin{equation}\tag{2.4.9}
 w:\mu_m\longrightarrow\Pic(Y,Z)^0.
\end{equation}
Passing to Lie algebras gives
\begin{equation}\tag{2.4.10}
 dw:\Lie\Gm\longrightarrow\Lie\Pic(Y;Z)^0,
\end{equation}
and hence an element
\begin{equation}\tag{2.4.11}
 dw(1)\in H^1(Y,\OO(-Z)).
\end{equation}

Let $\omega_{Y/S}$ be the relative dualizing sheaf and let $\varphi$ be a section of $\omega_{Y/S}(Z)$. In formal coordinates near $O$,
\[
 \varphi=f(x,y)\frac{dx\wedge dy}{dt}.
\]
The restriction of $\varphi$ to the generic fibre is paired with $H^1(Y_\eta,\OO(-Z_\eta))$ by Serre duality.

\paragraph{2.5. Proposition.}
With the notation above,
\begin{equation}\tag{2.5.1}
 I(f)=-(dw(1),\varphi).
\end{equation}

The right-hand side lies in $k[[t]]$, not merely in $k((t))$. Reducing modulo $t^n$, the local trace and the global Serre trace are compatible; the proposition is reduced to the following lemma.

\paragraph{2.6. Lemma.}
The morphism from local cohomology at $O$ to $H^1(Y_n,\OO(-Z_n))$ sends $v_n$ to $dw_n(1)$.

The image of $v_n$ is also the image of $1$ under (2.4.6). The assertion follows from the description of $dw_n$ via dual numbers over $k[[t]]/(t^n)$ and the natural exact sequence comparing the additive and multiplicative constructions.

\section{Proof of the theorem}

Let $S_0$ be a smooth curve over $k$, let $s\in S_0(k)$, and let $S$ be the completion at $s$. Let $\pi:Y_0\to S_0$ be proper and flat, purely of relative dimension one, with $Y$ its base change to $S$ and with a point $O\in Y_s$ such that $(Y/S,O)$ is as above. We suppose, for simplicity, that $Y_0$ is regular and that $\pi$ has geometrically connected fibres. Let $\omega_{Y_0/S_0}$ be the relative dualizing sheaf.

Let $\varphi$ be a rational section of $\omega_{Y_0/S_0}$, regular at $O$. In formal coordinates it can be written
\[
 \varphi=f(x,y)\frac{dx\wedge dy}{dt}.
\]

\paragraph{3.2. Proposition.}
If $k$ has characteristic $p>0$, then $I(f)$ is algebraic over the field $k(S_0)$ of rational functions on $S_0$.

After multiplying $\varphi$ by a rational function on $S_0$ vanishing sufficiently at $s$, and after shrinking $S_0$, we may assume that for a suitable relative Cartier divisor $Z_0$ disjoint from $O$, the form $\varphi$ is a section of $\omega$ on $Y_0-Z_0$. Let $\Pic(Y_0;Z_0)^0$ be the relative generalized Picard scheme. Via Serre duality, $\varphi$ corresponds to a section of the dual of the vector bundle
\[
 \Lie\Pic(Y_0;Z_0)^0=R^1\pi_*\OO(-Z_0).
\]
After base change to $S$ we have $dw:\OO\to\Lie\Pic(Y;Z)^0$, and Proposition 2.5 reduces Proposition 3.2 to the following key lemma.

\paragraph{3.3. Key Lemma.}
If $p>0$, then $dw$ is algebraic over $S_0$.

In characteristic $p$, the inclusion $\mu_p\subset\Gm$ induces an isomorphism on Lie algebras:
\[
 \Lie(\mu_p)\xrightarrow{\sim}\Lie(\Gm).
\]
Modulo $t^n$, one checks that $dw$ is the differential of the morphism
\[
 w:\mu_p\to\Pic(Y;Z)^0.
\]
It remains to use the rigidity of $\mu_p$: for every flat commutative group scheme $P$ over $S_0$, the functor $\Hom(\mu_p,P)$ is represented by an etale scheme over $S_0$. Applying this to $P=\Pic(Y_0;Z_0)^0$, the morphism $w$ is pulled back from a universal morphism over an etale cover of $S_0$. Hence its differential, and therefore $dw$, is algebraic.

Let $k\{x_1,\ldots,x_N\}$ be the henselization of the local ring of affine $N$-space at the origin. Its completion is $k[[x_1,\ldots,x_N]]$, and an element of the completion is algebraic over $k(x_1,\ldots,x_N)$ if and only if it lies in the henselization.

\subsection*{Proof of the theorem}

If $f$ is algebraic, there exists a surface $Y_0$ with a rational point $O$, an etale morphism $h:Y_0\to\A^2$ sending $O$ to the origin, and a section whose formal expansion at $O$ is $f$. Let $S_0=\A^1$ with coordinate $t$, and let $\pi'$ be the composite
\[
 Y_0\xrightarrow{h}\A^2\xrightarrow{xy}\A^1=S_0.
\]
The differential form $f\,h^*((dx\wedge dy)/dt)$ is a rational section of the relative dualizing sheaf. After compactifying and resolving singularities, one obtains a proper flat regular model to which Proposition 3.2 applies. Therefore $I(f)$ is algebraic.

\subsection*{Proof of the corollary}

Proceed by induction on $N$. The cases $N=0$ and $N=1$ are immediate, and $N=2$ is the theorem. For $N>2$, write
\[
 g=\sum g_{n,m}x_{N-1}^n x_N^m,
\qquad g_{n,m}\in k[[x_1,\ldots,x_{N-2}]],
\]
and put
\[
 I(g)=\sum g_{n,n}x_{N-1}^n.
\]
Then
\begin{equation}\tag{3.6.1}
 I_N(g)=I_{N-1}(I(g)).
\end{equation}
If $g$ is algebraic, the coefficients $g_{n,m}$ are algebraic over $k(x_1,\ldots,x_{N-2})$. Applying the theorem over the fraction field of the henselization in the first $N-2$ variables shows that $I(g)$ is algebraic, and induction proves the claim.

\subsection*{Degree estimates and remarks}

The proof gives an estimate. With the notation of the geometric construction, let $\alpha$ be the degree of the curve $S_0$ over the $t$-line, let $g$ be the genus of the normalization of the geometric generic fibre, and let $z=\max(0,z'-1)$ where $z'$ is the number of geometric points of the divisor $Z$. Then the degree $d$ of the equation satisfied by $I(f)$ is bounded by
\begin{equation}\tag{3.7.1}
 d<\alpha p^{g+z}.
\end{equation}
For a formal algebraic series with integer coefficients, this implies that for almost all primes $p$, the degree $d_p$ of the reduction of $I(f)$ modulo $p$ satisfies
\begin{equation}\tag{3.7.2}
 d_p<\alpha p^{g+z}.
\end{equation}
Using $\mu_{p^m}$ when $p^m=0$ on the base gives analogous information modulo powers of $p$, for almost all primes.

A direct proof for $N$ variables would be preferable. Analytically, for a convergent series $g$ one expects the formula
\begin{equation}\tag{1'}
 I(g)(t)=\int_{Z(t)}g\,\frac{dz_1\cdots dz_N}{dt},
\end{equation}
where
\[
 Z(t)=\{(x_1,\ldots,x_N):x_1\cdots x_N=t,
 |x_i|=r_i,
 \prod r_i=|t|\}.
\]
A direct proof of this kind would be a test case for a $p$-adic theory of vanishing cycles.

Furstenberg's proof in the rational case shows that if $g$ is one component of a column vector $h$ with entries in $k[[x_1,\ldots,x_N]]$ satisfying
\[
 h=A h^p
\]
for a square matrix $A$ with polynomial entries, then $I(g)$ has the same property and is algebraic. This type of equation recalls $F$-crystals in crystalline cohomology. Since $T_p\Pic$ is an avatar of $p$-adic $H^1$, the proof here and Furstenberg's proof are perhaps less distant from each other than they first appear.

\appendix
\section{Signs}

The treatment of signs in the duality formalism for coherent algebraic sheaves is delicate. We record the conventions used here. For a flat relative complete intersection morphism $f:X\to S$, purely of relative dimension $n$, the relative dualizing sheaf $\omega_{X/S}$ is related to the relative dualizing complex by
\begin{equation}\tag{S.1}
 K_{X/S}=\omega_{X/S}[n].
\end{equation}
When $X$ is smooth over $S$, $\omega_{X/S}=\Omega^n_{X/S}$. If $X$ is the fibre product over $S$ of morphisms $f_i:X_i\to S$ of relative dimensions $n_i$, then
\begin{equation}\tag{S.2}
 K_{X/S}=\boxtimes_i K_{X_i/S}.
\end{equation}
Equivalently, after choosing a total order on the index set,
\begin{equation}\tag{S.3}
 \omega_{X/S}=\bigotimes_i\omega_{X_i/S},
\end{equation}
with the sign change $(-1)^N$ when the order is changed, where $N$ is the sum of $n_in_j$ over inverted pairs.

For $f:X\to S$ proper as above, the trace is
\begin{equation}\tag{S.4}
 \Tr:Rf_*K_{X/S}\to\OO_S.
\end{equation}
It induces
\begin{equation}\tag{S.5}
 \Tr:R^nf_*\omega_{X/S}\to\OO_S.
\end{equation}
For a section $s$ of $f$ there is a local trace
\begin{equation}\tag{S.6}
 \Tr:Rs^!K_{X/S}\to\OO_S.
\end{equation}
The choices are fixed by requiring compatibility between global and local traces and by normalizing, for a smooth curve $X$ over an algebraically closed field and a point $x\in X(k)$, the local trace
\begin{equation}\tag{S.7}
 H^1_{\{x\}}(X,\Omega^1[1])\to k
\end{equation}
to be the residue map $\Res_x$.

With these conventions, for a smooth complete curve $X$ and $x\in X(k)$, the boundary map
\[
 H^0(X-\{x\},\Omega^1)\to H^1_{\{x\}}(X,\Omega^1)
\]
composed with the trace is $-\Res_x$. For products, if $X=X_1\times X_2$ and $u_i\in H^{n_i}(X_i,\omega_{X_i})$, then
\[
 \Tr(\operatorname{pr}_1^*u_1\smile\operatorname{pr}_2^*u_2)
 =(-1)^{n_1n_2}\Tr(u_1)\Tr(u_2).
\]
Over $k=\C$, comparison with the transcendental trace identifies the algebraic trace with integration over $X(\C)$, with the standard normalization by $(2\pi i)^{-n}$. For projective space $\mathbf P^n$ with homogeneous coordinates $T_0,\ldots,T_n$ and $t_i=T_i/T_0$, the Cech class of
\[
 \frac{dt_1\wedge\cdots\wedge dt_n}{t_1\cdots t_n}
\]
has trace $(-1)^{n(n+1)/2}$ with these conventions.

\begin{thebibliography}{9}
\bibitem{Furstenberg} H. Furstenberg, Algebraic functions over finite fields, \emph{J. Algebra} \textbf{7} (1967), 271--277.
\bibitem{RaynaudPic} M. Raynaud, Specialisation du foncteur de Picard, \emph{Publ. Math. IHES} \textbf{38} (1970), 27--76.
\bibitem{RaynaudHensel} M. Raynaud, \emph{Anneaux locaux henseliens}, Lecture Notes in Mathematics 169, Springer, 1970.
\end{thebibliography}

\end{document}
